\chapter{Kết luận và hướng phát triển}
\label{ch:ketluan}

\begin{muctieu}
Chương này tổng kết nguyên mẫu đã xây dựng, nêu giới hạn và các bước phát triển tiếp theo.
\end{muctieu}

\section{Kết luận}
Đồ án xây dựng nguyên mẫu trợ lý tuyển dụng gồm giao diện quản lý, API backend, lưu dữ liệu, đọc CV và tính điểm so khớp. Hệ thống kết hợp các quy tắc trích xuất, độ tương đồng văn bản và mô hình Random Forest. Trên tập đánh giá hiện có, macro F1 của Random Forest cao hơn baseline lớp đa số, trong khi accuracy thấp hơn. Do dữ liệu đánh giá có thành phần tổng hợp và nhãn hỗ trợ bằng mô hình, chưa thể kết luận hiệu quả trên quy trình tuyển dụng thực tế.

\section{Bài học rút ra}
Chất lượng dữ liệu, nhãn và ngữ cảnh hồ sơ ảnh hưởng trực tiếp đến kết quả. Chỉ số tổng hợp không thay thế việc kiểm tra lỗi từng lớp và từng trường trích xuất. Hệ thống AI trong tuyển dụng cần cho phép con người rà soát, giải thích được tín hiệu chính và tránh tự động hóa quyết định bất lợi cho ứng viên.

\section{Hướng phát triển}
\begin{itemize}[leftmargin=1.6em,itemsep=2pt]
  \item Thu thập dữ liệu có sự đồng ý và xây dựng tập nhãn độc lập, có quy trình kiểm tra chất lượng.
  \item Đánh giá riêng độ chính xác trích xuất email, kỹ năng, học vấn và khoảng thời gian kinh nghiệm.
  \item Cải thiện phân đoạn bố cục và nhận diện ngữ cảnh thời gian để giảm nhầm lẫn giữa học vấn và kinh nghiệm.
  \item Kiểm tra thiên lệch theo nhóm phù hợp với quy định và dữ liệu có thể sử dụng; lưu giải thích các thành phần điểm.
  \item Hoàn thiện kiểm thử tích hợp, bảo vệ dữ liệu cá nhân và đánh giá với người dùng tuyển dụng.
\end{itemize}
